% =============================================================================
% APPENDIX TS: LIT-SCHELLING: Thomas Schelling Research Integration
% =============================================================================
% Module: BCM2_LIT_SCHELLING (Literature Integration)
% Category: LIT
% Version: 1.0 (January 2026)
% Status: Final
% Language: English
%
% This appendix integrates Thomas Schelling's foundational contributions to
% strategic interaction, focal points, segregation dynamics, and self-command
% into the Evidence-Based Framework (EBF).
%
% =============================================================================

% =============================================================================
% CROSS-REFERENCE STRUCTURE
% =============================================================================
%
% CHAPTER LINKAGE (Required):
% - Primary Chapter: Chapter 5 (Complementarity)
% - Secondary Chapters: Chapter 9 (Context), Chapter 12 (Willingness),
%                       Chapter 14 (Behavioral Change)
%
% DEPENDENCIES (what this appendix REQUIRES):
% - Appendix CA (LIT-CAMERER): Game theory foundations
% - Appendix K (LIT-FEHR): Social preferences experiments
% - Appendix AG (DOMAIN-COMPLEX): Complexity economics
%
% DEPENDENTS (what REQUIRES this appendix):
% - Appendix AV (CORE-READY): Self-command theory
% - Appendix AW (CORE-STAGE): Commitment devices
% - Appendix HF (PREDICT-HABIT): Self-control and addiction
%
% =============================================================================

\begin{tcolorbox}[colback=purple!5!white, colframe=purple!75!black,
    title=Chapter Linkage]

\textbf{Primary Chapter:} Chapter 5: Complementarity
\begin{itemize}[nosep]
\item Section 5.3: Strategic Interaction $\leftrightarrow$ This Appendix Section TS.2
\item Section 5.5: Coordination Problems $\leftrightarrow$ This Appendix Section TS.3
\end{itemize}

\textbf{Secondary Chapters:}
\begin{itemize}[nosep]
\item Chapter 9: Context --- focal points as context-dependent coordination
\item Chapter 12: Willingness --- self-command and commitment devices
\item Chapter 14: Behavioral Change --- tipping points and critical mass
\end{itemize}

\textbf{Reading Guide:}
\begin{itemize}[nosep]
\item For strategic interaction: Read this appendix first
\item For game theory foundations: See Appendix CA (LIT-CAMERER)
\item For complexity applications: See Appendix AG (DOMAIN-COMPLEX)
\end{itemize}

\end{tcolorbox}

\chapter{LIT-SCHELLING: Thomas Schelling Research Integration}
\label{app:lit-schelling}


\begin{tcolorbox}[colback=blue!5!white, colframe=blue!75!black,
    title=Quick Reference: Appendix TS]

\textbf{Appendix:} TS (LIT)

\textbf{Core Question:} What are the key findings in this domain?

\textbf{Key Concepts:}
\begin{itemize}[nosep]
\item Framework integration
\item Parameter validation
\item Cross-references
\end{itemize}

\textbf{Cross-References:} See related appendices below
\end{tcolorbox}

\begin{abstract}
Thomas Schelling (Nobel Prize 2005) pioneered the analysis of strategic interaction
beyond zero-sum games, introducing focal points, commitment strategies, and tipping
point dynamics. This appendix integrates his 50+ years of research on conflict
resolution, segregation dynamics, self-command, and micromotives-to-macrobehavior
aggregation into EBF. Schelling's work provides the theoretical foundation for
understanding how individual strategic choices aggregate into collective patterns,
and how context shapes coordination without explicit communication.
\end{abstract}

\begin{tcolorbox}[colback=blue!5!white, colframe=blue!75!black,
    title=Quick Reference: LIT-SCHELLING]

\textbf{Core Question:} How do individuals coordinate on shared expectations
(focal points), and how do individual micromotives aggregate into emergent
macrobehavior---including segregation, tipping, and self-command failures?

\textbf{Key Formula (Focal Point Coordination):}
\begin{equation}
P(x^* | \Psi) = \frac{\text{Salience}(x^*, \Psi)}{\sum_{x'} \text{Salience}(x', \Psi)}
\label{eq:ts-focal}
\end{equation}

\textbf{Key Concepts:}
\begin{itemize}[nosep]
\item \textbf{Focal Points (Schelling Points)}: Context-dependent coordination solutions
\item \textbf{Tipping Points}: Critical thresholds where individual choices cascade
\item \textbf{Commitment}: Strategic self-binding to influence others' expectations
\item \textbf{Self-Command}: The internal strategic problem of multiple selves
\end{itemize}

\textbf{Cross-References:} LIT-CAMERER (CA), LIT-RABIN (RA), CORE-READY (AV), DOMAIN-COMPLEX (AG)
\end{tcolorbox}

% -----------------------------------------------------------------------------
% APPENDIX SCOPE BOX
% -----------------------------------------------------------------------------

\begin{tcolorbox}[
    colback=orange!5!white,
    colframe=orange!75!black,
    title={\textbf{Appendix Scope: Ziel / In-Scope / Out-of-Scope / Constraints}},
    fonttitle=\bfseries
]

\textbf{Ziel (Objective):}
\begin{itemize}[nosep]
    \item How do contextual salience and shared expectations enable coordination
in the absence of communication, and how do individual micromotives aggregate
into emergent macrobehavior---including patterns that no individual intended?
\end{itemize}

\textbf{In-Scope:}
\begin{itemize}[nosep]
    \item The Fundamental Question
    \item Core Theory: The Schelling Research Program
    \item Axioms: The Schelling Framework
    \item Integration with EBF
    \item Worked Example
\end{itemize}

\textbf{Out-of-Scope (delegiert an andere Appendices):}
\begin{itemize}[nosep]
    \item LIT-CAMERER $\rightarrow$ Appendix CA
    \item LIT-FEHR $\rightarrow$ Appendix K
    \item DOMAIN-COMPLEX $\rightarrow$ Appendix AG
    \item CORE-READY $\rightarrow$ Appendix AV
    \item CORE-STAGE $\rightarrow$ Appendix AW
\end{itemize}

\textbf{Constraints (Anwendungsgrenzen):}
\begin{itemize}[nosep]
    \item [TODO: Define when this appendix does NOT apply]
    \item [TODO: Define required prerequisites]
    \item [TODO: Define known limitations]
\end{itemize}

\textbf{Lieferobjekte (Deliverables):}
\begin{itemize}[nosep]
    \item 3 Table(s)
    \item 17 Equation(s)
    \item 7 Axiom(s)
    \item Worked Example(s)
\end{itemize}

\end{tcolorbox}




%%%%%%%%%%%%%%%%%%%%%%%%%%%%%%%%%%%%%%%%%%%%%%%%%%%%%%%%%%%%%%%%%%%%%%%%%%%%%%%
\section{The Fundamental Question}
\label{sec:ts-fundamental}
%%%%%%%%%%%%%%%%%%%%%%%%%%%%%%%%%%%%%%%%%%%%%%%%%%%%%%%%%%%%%%%%%%%%%%%%%%%%%%%

\subsection{Why This Matters}

Before Schelling's work, game theory focused primarily on zero-sum games with
clear equilibria. Schelling revolutionized the field by asking: How do people
coordinate when multiple equilibria exist? His answer---focal points---revealed
that context, culture, and shared expectations determine which equilibrium emerges.

His work matters for EBF because:
\begin{enumerate}
\item \textbf{Focal Points}: Shows how $\Psi$-context determines coordination
      without explicit communication
\item \textbf{Tipping Points}: Explains how individual preferences aggregate into
      segregation and critical mass dynamics
\item \textbf{Self-Command}: Provides the theoretical foundation for CORE-READY
      and the internal struggle for self-control
\item \textbf{Commitment}: Shows how strategic self-binding can overcome
      present bias and temptation
\end{enumerate}

\subsection{The Core Problem}

Classical game theory assumed common knowledge and unique equilibria. Schelling
recognized that:

\begin{quote}
\textit{``People can often concert their intentions or expectations with others
if each knows that the other is trying to do the same.''}
--- Schelling (1960)
\end{quote}

\begin{tcolorbox}[colback=red!5!white, colframe=red!75!black,
    title=The Fundamental Question]
\textbf{How do contextual salience and shared expectations enable coordination
in the absence of communication, and how do individual micromotives aggregate
into emergent macrobehavior---including patterns that no individual intended?}
\end{tcolorbox}


%%%%%%%%%%%%%%%%%%%%%%%%%%%%%%%%%%%%%%%%%%%%%%%%%%%%%%%%%%%%%%%%%%%%%%%%%%%%%%%
\section{Core Theory: The Schelling Research Program}
\label{sec:ts-theory}
%%%%%%%%%%%%%%%%%%%%%%%%%%%%%%%%%%%%%%%%%%%%%%%%%%%%%%%%%%%%%%%%%%%%%%%%%%%%%%%

\subsection{Focal Points and Tacit Coordination (1960)}

Schelling's most famous contribution is the concept of \textit{focal points}---
solutions to coordination games that are salient because of contextual features,
not strategic dominance.

\textbf{The Classic Example:}
``Name a place to meet in New York City tomorrow.'' Most people say ``Grand Central
Station at noon''---not because it's strategically optimal, but because it's
\textit{salient} in the shared cultural context.

\begin{equation}
x^* = \arg\max_x \text{Salience}(x | \Psi_{shared})
\label{eq:ts-focal-max}
\end{equation}

\textbf{Key Insight:} Focal points are context-dependent. Different cultures,
different experiences, different $\Psi$ produce different focal solutions.

\textbf{EBF Integration:} Focal points demonstrate that CORE-WHEN ($\Psi$)
determines equilibrium selection in coordination games. The ``right'' answer
depends on shared context, not individual optimization.

\subsection{Segregation and Tipping Points (1971, 1978)}

Schelling's segregation model demonstrated that mild individual preferences for
similar neighbors can produce complete segregation---a macrobehavior far exceeding
individual micromotives.

\textbf{The Segregation Dynamic:}
\begin{equation}
\text{Move}_i = \mathbf{1}\left[\frac{\text{Similar}_i}{\text{Total}_i} < \tau_i\right]
\label{eq:ts-segregation}
\end{equation}

where $\tau_i \approx 0.33$ (individuals tolerate up to 2/3 different neighbors).

\textbf{Key Finding:} Even when no individual \textit{prefers} segregation, the
aggregate outcome can be complete segregation. Individual tolerance does not
produce collective integration.

\begin{equation}
\text{Aggregate Segregation} > \max_i \{\text{Preferred Segregation}_i\}
\label{eq:ts-emergent}
\end{equation}

\textbf{Tipping Point Dynamics:}
\begin{equation}
\frac{dN_{minority}}{dt} = -\alpha \cdot (N_{minority} - N_{critical}) \quad \text{for } N_{minority} < N_{critical}
\label{eq:ts-tipping}
\end{equation}

Once the minority population falls below a critical threshold, the neighborhood
``tips'' toward homogeneity.

\subsection{Self-Command and the Multiple-Self Problem (1980, 1984)}

Schelling extended strategic analysis inward, treating the individual as a
collection of competing ``selves'' across time.

\textbf{The Self-Command Problem:}
\begin{equation}
\max_{x_t} U_0(x_0, x_1, ..., x_T) \quad \text{subject to} \quad x_t = \arg\max U_t(x_t, ..., x_T)
\end{equation}

The ``present self'' wants to commit the ``future self'' to certain behaviors,
but lacks direct control.

\textbf{Commitment Devices:}
\begin{itemize}
\item \textbf{Physical}: Removing temptation (no cookies in the house)
\item \textbf{Social}: Public announcements (telling friends about diet)
\item \textbf{Financial}: Penalties for deviation (gym membership contracts)
\item \textbf{Temporal}: Cooling-off periods (waiting periods for purchases)
\end{itemize}

\textbf{EBF Integration:} Schelling's self-command framework provides the foundation
for CORE-READY (AV) and the WAX threshold mechanism. The struggle is not just
between individual and environment, but between \textit{selves across time}.

\subsection{Strategic Commitment and Credibility (1956, 1960)}

Schelling analyzed how limiting one's own options can paradoxically increase
bargaining power.

\textbf{The Commitment Paradox:}
\begin{equation}
\text{Bargaining Power} = f(\text{Credibility of Threat}) = g(\text{Cost of Retreat})
\label{eq:ts-commitment}
\end{equation}

By making retreat costly or impossible, a negotiator can credibly commit to a
position and extract concessions.

\textbf{Applications:}
\begin{itemize}
\item Burning bridges (making retreat impossible)
\item Reputation staking (making deviation costly)
\item Delegation (removing decision authority)
\item Automation (removing human discretion)
\end{itemize}


%%%%%%%%%%%%%%%%%%%%%%%%%%%%%%%%%%%%%%%%%%%%%%%%%%%%%%%%%%%%%%%%%%%%%%%%%%%%%%%
\section{Axioms: The Schelling Framework}
\label{sec:ts-axioms}
%%%%%%%%%%%%%%%%%%%%%%%%%%%%%%%%%%%%%%%%%%%%%%%%%%%%%%%%%%%%%%%%%%%%%%%%%%%%%%%

\begin{tcolorbox}[colback=gray!5!white, colframe=gray!75!black,
    title=Axiom TS-1: Focal Point Salience]
\textbf{Statement:} In coordination games with multiple equilibria, players
converge on solutions that are salient in the shared context $\Psi$, not
strategically dominant.

\begin{equation}
P(x^* | \text{coordination game}) = f(\text{Salience}(x^*, \Psi_{shared}))
\end{equation}

\textbf{Interpretation:} Context determines coordination outcomes. Different
$\Psi$ produces different focal points, even for identical strategic structures.

\textbf{Implication:} EBF's CORE-WHEN ($\Psi$) determines equilibrium selection,
not just parameter values.
\end{tcolorbox}

\begin{tcolorbox}[colback=gray!5!white, colframe=gray!75!black,
    title=Axiom TS-2: Emergent Macrobehavior]
\textbf{Statement:} Aggregate outcomes can systematically exceed individual
intentions---small micromotives can produce large macrobehavior patterns.

\begin{equation}
\text{Outcome}_{aggregate} \neq f(\text{Outcome}_{intended})
\end{equation}

\textbf{Interpretation:} The whole is not the sum of parts. Individual tolerance
can produce collective intolerance; individual rationality can produce collective
irrationality.

\textbf{Implication:} EBF's CORE-WHO (AAA) must account for emergent level
properties that exceed individual-level specifications.
\end{tcolorbox}

\begin{tcolorbox}[colback=gray!5!white, colframe=gray!75!black,
    title=Axiom TS-3: Tipping Point Dynamics]
\textbf{Statement:} Many social phenomena exhibit critical thresholds below
which dynamics reverse and outcomes cascade.

\begin{equation}
\frac{dN}{dt} = \alpha \cdot (N - N_{crit}) \quad \text{with sign change at } N_{crit}
\end{equation}

\textbf{Interpretation:} Small changes near the tipping point produce
disproportionate effects. Path dependence emerges from threshold dynamics.

\textbf{Implication:} EBF interventions must consider tipping dynamics---
incremental changes may fail, while threshold-crossing changes succeed.
\end{tcolorbox}

\begin{tcolorbox}[colback=gray!5!white, colframe=gray!75!black,
    title=Axiom TS-4: Self-Command as Strategic Problem]
\textbf{Statement:} Self-control is a strategic interaction between temporal
selves with conflicting preferences.

\begin{equation}
\max_{c_0} U_0(c_0, c_1^*) \quad \text{where } c_1^* = \arg\max_{c_1} U_1(c_1)
\end{equation}

\textbf{Interpretation:} The present self cannot directly control future selves,
only influence them through commitment devices and environmental structuring.

\textbf{Implication:} CORE-READY (AV) must model the WAX threshold as a
commitment problem, not just an awareness problem.
\end{tcolorbox}

\begin{tcolorbox}[colback=gray!5!white, colframe=gray!75!black,
    title=Axiom TS-5: Commitment Credibility]
\textbf{Statement:} Commitment power increases with the cost of deviation---
constraining one's own options can increase bargaining power.

\begin{equation}
\text{Credibility}(\text{threat}) = f(\text{Cost}(\text{not executing threat}))
\end{equation}

\textbf{Interpretation:} Strategic self-limitation (burning bridges, staking
reputation) can be rational when it influences others' expectations.

\textbf{Implication:} EBF commitment devices work by making deviation costly,
not by making the committed action more attractive.
\end{tcolorbox}

\begin{tcolorbox}[colback=gray!5!white, colframe=gray!75!black,
    title=Axiom TS-6: Statistical Lives vs. Identified Lives]
\textbf{Statement:} People value identified lives (specific individuals at risk)
more than statistical lives (anonymous members of a population).

\begin{equation}
V(\text{save identified life}) > V(\text{save statistical life})
\end{equation}

\textbf{Interpretation:} This ``identifiable victim effect'' explains why rescue
operations receive more resources than prevention programs.

\textbf{Implication:} EBF interventions should consider whether targets are
identified or statistical when designing communication and motivation.
\end{tcolorbox}


%%%%%%%%%%%%%%%%%%%%%%%%%%%%%%%%%%%%%%%%%%%%%%%%%%%%%%%%%%%%%%%%%%%%%%%%%%%%%%%

%%%%%%%%%%%%%%%%%%%%%%%%%%%%%%%%%%%%%%%%%%%%%%%%%%%%%%%%%%%%%%%%%%%%%%%%%%%%%%%
\section{Key Results}
\label{sec:results}
%%%%%%%%%%%%%%%%%%%%%%%%%%%%%%%%%%%%%%%%%%%%%%%%%%%%%%%%%%%%%%%%%%%%%%%%%%%%%%%

The analysis yields the following key results:

\begin{enumerate}
    \item \textbf{Result 1:} [TODO: Describe main finding]
    \item \textbf{Result 2:} [TODO: Describe secondary finding]
    \item \textbf{Result 3:} [TODO: Describe implication]
\end{enumerate}

\section{Integration with EBF}
\label{sec:ts-integration}
%%%%%%%%%%%%%%%%%%%%%%%%%%%%%%%%%%%%%%%%%%%%%%%%%%%%%%%%%%%%%%%%%%%%%%%%%%%%%%%

\subsection{Connection to CORE-WHEN (V)}

Schelling's focal point theory directly supports EBF's context framework:

\begin{table}[htbp]
\centering
\caption{Schelling Concepts $\leftrightarrow$ EBF Context}
\label{tab:ts-psi}
\renewcommand{\arraystretch}{1.3}
\begin{tabular}{@{}lll@{}}
\toprule
\textbf{Schelling Concept} & \textbf{EBF Component} & \textbf{Application} \\
\midrule
Focal Point Salience & $\Psi_{cultural}$ & Equilibrium selection \\
Shared Expectations & $\Psi_{social}$ & Coordination norms \\
Physical Environment & $\Psi_{spatial}$ & Choice architecture \\
Temporal Context & $\Psi_{temporal}$ & Self-command timing \\
\bottomrule
\end{tabular}
\end{table}

\subsection{Connection to CORE-READY (AV)}

Schelling's self-command work provides the theoretical foundation for the
WAX threshold:

\begin{equation}
WAX = f(\text{Self-Control Capacity}, \text{Commitment Devices}, \Psi_{temptation})
\label{eq:ts-wax}
\end{equation}

The ``intimate contest for self-command'' maps directly to the tension between
$WAX$ and $\theta$ (action threshold):
\begin{itemize}
\item \textbf{Self-Control Capacity}: Baseline WAX level
\item \textbf{Commitment Devices}: Mechanisms to raise effective WAX
\item \textbf{Temptation Context}: $\Psi$ factors that lower WAX
\end{itemize}

\subsection{Connection to CORE-WHO (AAA)}

Schelling's micromotives-macrobehavior analysis supports EBF's level structure:

\begin{equation}
U^{Level}_{aggregate} \neq \sum_i U^{Individual}_i
\label{eq:ts-levels}
\end{equation}

Emergent properties at higher levels (segregation, tipping) arise from
interactions at lower levels, creating genuine level-specific utility
that cannot be reduced to individual preferences.

\begin{tcolorbox}[colback=yellow!5!white, colframe=yellow!75!black,
    title=Integration: LIT-SCHELLING $\leftrightarrow$ EBF CORE]
\textbf{What LIT-SCHELLING provides:}
\begin{itemize}[nosep]
\item Focal point theory for $\Psi$-dependent coordination
\item Tipping point dynamics for behavioral change thresholds
\item Self-command framework for CORE-READY (AV)
\item Micromotives-macrobehavior aggregation for CORE-WHO (AAA)
\end{itemize}

\textbf{What it requires:}
\begin{itemize}[nosep]
\item CORE-WHEN for context specification of focal points
\item CORE-READY for self-command and commitment operationalization
\item CORE-WHO for level aggregation mechanisms
\item DOMAIN-COMPLEX for tipping point dynamics
\end{itemize}
\end{tcolorbox}


%%%%%%%%%%%%%%%%%%%%%%%%%%%%%%%%%%%%%%%%%%%%%%%%%%%%%%%%%%%%%%%%%%%%%%%%%%%%%%%
\section{Worked Example}
\label{sec:ts-example}
%%%%%%%%%%%%%%%%%%%%%%%%%%%%%%%%%%%%%%%%%%%%%%%%%%%%%%%%%%%%%%%%%%%%%%%%%%%%%%%

\subsection{Segregation Dynamics and EBF Levels}

\paragraph{Setup.} Consider a neighborhood with 100 residents, 50\% Group A
and 50\% Group B. Each individual has tolerance threshold $\tau = 0.33$
(comfortable if at least 1/3 of neighbors are similar).

\paragraph{Application.} Using Schelling's segregation model:

\textbf{Step 1: Individual Preferences}
\begin{equation}
\text{Satisfied}_i = \mathbf{1}\left[\frac{N_{similar,i}}{N_{total,i}} \geq 0.33\right]
\end{equation}

\textbf{Step 2: Initial State}
Random distribution: most individuals satisfied (average 50\% similar).

\textbf{Step 3: Perturbation}
One dissatisfied individual moves, changing neighbor composition for others.

\textbf{Step 4: Cascade}
Moving individuals create new dissatisfied individuals, triggering further moves.

\textbf{Step 5: Equilibrium}
System reaches stable state with $>80\%$ segregation---far exceeding individual
preference for 33\% similar.

\paragraph{Result.} Individual tolerance ($\tau = 0.33$) produces collective
segregation ($>80\%$). The emergent macrobehavior exceeds any individual's
preferred outcome.

\paragraph{EBF Interpretation.}
\begin{equation}
U^{Neighborhood}_{actual} < U^{Neighborhood}_{preferred} \quad \text{(collective welfare loss)}
\end{equation}

This demonstrates Axiom TS-2: aggregate outcomes systematically exceed individual
intentions. CORE-WHO must account for emergent properties at collective levels.


%%%%%%%%%%%%%%%%%%%%%%%%%%%%%%%%%%%%%%%%%%%%%%%%%%%%%%%%%%%%%%%%%%%%%%%%%%%%%%%
\section{Implications for Practice}
\label{sec:ts-practice}
%%%%%%%%%%%%%%%%%%%%%%%%%%%%%%%%%%%%%%%%%%%%%%%%%%%%%%%%%%%%%%%%%%%%%%%%%%%%%%%

\begin{tcolorbox}[colback=green!5!white, colframe=green!75!black,
    title=Practical Implications]
\begin{enumerate}
\item \textbf{Focal Point Design}: Create salient coordination points through
      choice architecture and environmental design
\item \textbf{Tipping Interventions}: Focus resources on crossing critical
      thresholds rather than incremental changes
\item \textbf{Commitment Devices}: Design self-binding mechanisms that raise
      the cost of deviation (not the benefit of compliance)
\item \textbf{Aggregate Awareness}: Monitor emergent macrobehavior, not just
      individual preferences---tolerant individuals can produce intolerant systems
\item \textbf{Statistical vs. Identified}: Frame interventions around identified
      beneficiaries when possible to increase motivation
\end{enumerate}
\end{tcolorbox}


%%%%%%%%%%%%%%%%%%%%%%%%%%%%%%%%%%%%%%%%%%%%%%%%%%%%%%%%%%%%%%%%%%%%%%%%%%%%%%%
\section{Summary}
\label{sec:ts-summary}
%%%%%%%%%%%%%%%%%%%%%%%%%%%%%%%%%%%%%%%%%%%%%%%%%%%%%%%%%%%%%%%%%%%%%%%%%%%%%%%

\begin{tcolorbox}[colback=green!10!white, colframe=green!75!black,
    title=\textbf{Appendix TS: Summary}]

\textbf{Core Question:} How do context and strategic interaction produce
coordination, commitment, and emergent collective patterns?

\textbf{Main Contribution:}
\begin{itemize}[nosep]
    \item Introduced focal points as context-dependent coordination solutions
    \item Demonstrated tipping point dynamics in segregation and social change
    \item Developed self-command theory for intertemporal self-control
    \item Analyzed strategic commitment and credibility mechanisms
    \item Bridged micromotives and macrobehavior aggregation
\end{itemize}

\textbf{Key Parameters:}
\begin{itemize}[nosep]
    \item Segregation threshold: $\tau \approx 0.33$ produces $>80\%$ segregation
    \item Focal point convergence: 70-90\% in experimental coordination games
    \item Tipping point: typically 15-30\% of population for adoption cascades
\end{itemize}

\textbf{Integration:} This appendix connects to CORE-WHEN (focal points),
CORE-READY (self-command), CORE-WHO (micromotives-macrobehavior), and
DOMAIN-COMPLEX (tipping dynamics).

\end{tcolorbox}


%%%%%%%%%%%%%%%%%%%%%%%%%%%%%%%%%%%%%%%%%%%%%%%%%%%%%%%%%%%%%%%%%%%%%%%%%%%%%%%
\section{Glossary of Symbols}
\label{sec:ts-glossary}
%%%%%%%%%%%%%%%%%%%%%%%%%%%%%%%%%%%%%%%%%%%%%%%%%%%%%%%%%%%%%%%%%%%%%%%%%%%%%%%

\begin{tcolorbox}[colback=blue!5!white, colframe=blue!75!black]
\textbf{Central Reference:} For complete symbol definitions, see
\textbf{Appendix G} (Glossary, \S\ref{app:glossary}).

This section lists only symbols introduced or specialized in this appendix.
\end{tcolorbox}

\begin{table}[htbp]
\centering
\caption{Symbols Introduced in Appendix TS}
\label{tab:ts-symbols}
\renewcommand{\arraystretch}{1.3}
\begin{tabular}{@{}clp{6cm}c@{}}
\toprule
\textbf{Symbol} & \textbf{Name} & \textbf{Definition} & \textbf{Also in G?} \\
\midrule
$x^*$ & Focal Point & Context-salient equilibrium & NEW \\
$\tau$ & Tolerance Threshold & Minimum similar neighbors & NEW \\
$N_{crit}$ & Tipping Point & Critical population threshold & NEW \\
$\Psi_{shared}$ & Shared Context & Common knowledge context & \checkmark \\
\bottomrule
\end{tabular}
\end{table}


%%%%%%%%%%%%%%%%%%%%%%%%%%%%%%%%%%%%%%%%%%%%%%%%%%%%%%%%%%%%%%%%%%%%%%%%%%%%%%%
\section{Critical Foundations}
\label{sec:ts-foundations}
%%%%%%%%%%%%%%%%%%%%%%%%%%%%%%%%%%%%%%%%%%%%%%%%%%%%%%%%%%%%%%%%%%%%%%%%%%%%%%%

\subsection{Objection 1: Model Simplicity}

\textbf{Objection:} Schelling's segregation model uses unrealistic assumptions
(binary groups, uniform preferences, grid topology) that limit generalizability.

\textbf{Response:} The power of Schelling's models lies precisely in their
simplicity---they demonstrate that even mild, uniform preferences can produce
extreme segregation. Extensions with heterogeneous preferences, continuous
variables, and network topologies confirm the basic mechanism.

\subsection{Objection 2: Focal Point Arbitrariness}

\textbf{Objection:} If focal points are culturally determined, they provide
no predictive power---any outcome can be rationalized as a focal point.

\textbf{Response:} Focal points are predictable within cultural contexts.
Experimental evidence shows high convergence rates (70-90\%) on focal solutions,
demonstrating that salience is not arbitrary but structured by shared experience.


%%%%%%%%%%%%%%%%%%%%%%%%%%%%%%%%%%%%%%%%%%%%%%%%%%%%%%%%%%%%%%%%%%%%%%%%%%%%%%%
\section{Open Issues}
\label{sec:ts-open}
%%%%%%%%%%%%%%%%%%%%%%%%%%%%%%%%%%%%%%%%%%%%%%%%%%%%%%%%%%%%%%%%%%%%%%%%%%%%%%%

\begin{table}[htbp]
\centering
\caption{Research Roadmap for LIT-SCHELLING Integration}
\label{tab:ts-roadmap}
\renewcommand{\arraystretch}{1.3}
\begin{tabular}{@{}clcl@{}}
\toprule
\textbf{\#} & \textbf{Issue} & \textbf{Priority} & \textbf{Status} \\
\midrule
1 & Focal point prediction across cultures & HIGH & PARTIAL \\
2 & Tipping point identification methods & HIGH & OPEN \\
3 & Self-command commitment device calibration & MEDIUM & PARTIAL \\
4 & Cross-level emergence formalization & MEDIUM & OPEN \\
\bottomrule
\end{tabular}
\end{table}


%%%%%%%%%%%%%%%%%%%%%%%%%%%%%%%%%%%%%%%%%%%%%%%%%%%%%%%%%%%%%%%%%%%%%%%%%%%%%%%
\section{References}
\label{sec:ts-references}
%%%%%%%%%%%%%%%%%%%%%%%%%%%%%%%%%%%%%%%%%%%%%%%%%%%%%%%%%%%%%%%%%%%%%%%%%%%%%%%

\begin{tcolorbox}[colback=blue!5!white, colframe=blue!75!black]
\textbf{Master Bibliography:} For complete reference details, see
\texttt{bcm2\_master\_references.bib} in the repository.
\nocite{bcm_master}
\end{tcolorbox}

\subsection{Key References for This Appendix}

\begin{itemize}[nosep]
    \item \textbf{Focal Points:} Schelling (1960, 1958)
    \item \textbf{Segregation:} Schelling (1971, 1969, 1978)
    \item \textbf{Self-Command:} Schelling (1980, 1984, 1985, 1992)
    \item \textbf{Commitment:} Schelling (1956, 1960, 1966)
    \item \textbf{Climate/Policy:} Schelling (1995, 2002, 2006)
\end{itemize}


%%%%%%%%%%%%%%%%%%%%%%%%%%%%%%%%%%%%%%%%%%%%%%%%%%%%%%%%%%%%%%%%%%%%%%%%%%%%%%%
% CROSS-REFERENCE FOOTER
%%%%%%%%%%%%%%%%%%%%%%%%%%%%%%%%%%%%%%%%%%%%%%%%%%%%%%%%%%%%%%%%%%%%%%%%%%%%%%%

\vspace{1em}
\hrule
\vspace{0.5em}

\noindent\textit{Cross-references:}
Appendix G (Glossary),
Appendix CA (LIT-CAMERER),
Appendix RA (LIT-RABIN),
Appendix AV (CORE-READY),
Appendix AG (DOMAIN-COMPLEX),
Appendix HF (PREDICT-HABIT).

\noindent\textit{Master Bibliography:} \texttt{bcm2\_master\_references.bib}

% =============================================================================
% END OF APPENDIX TS
% =============================================================================
